\documentclass[a4paper, 12pt]{article}
\usepackage[english]{babel}
\usepackage[utf8]{inputenc}
\usepackage[T1]{fontenc}
\usepackage{geometry}
\usepackage{parskip}
\usepackage{titlesec}
\usepackage{hyperref}
\usepackage{graphicx}
\usepackage{listings}
\usepackage{xcolor}

\geometry{margin=2.5cm}

\titleformat{\section}{\normalfont\Large\bfseries}{\thesection}{1em}{}
\titleformat{\subsection}{\normalfont\large\bfseries}{\thesubsection}{1em}{}

% Code listing configuration
\lstset{
    basicstyle=\ttfamily\small,
    breaklines=true,
    frame=single,
    backgroundcolor=\color{gray!10}
}

\title{StreamlitDRV: A Streamlit-Based Data Visualization Tool\\ \vspace{1.5em}
Project Report}
\author{Wojciech Bartoszek \\ Bartosz Gacek \\ Jarosław Kołdun}
\date{2025}

\begin{document}

\maketitle

\section{Introduction}

StreamlitDRV is an interactive web application for visualizing data through dimensionality reduction methods. The project is built on the Streamlit framework, which enables rapid development of analytical web applications. The primary goal of the application is to allow users to explore their own datasets and compare various dimensionality reduction techniques in an intuitive way.

The application addresses the need to visualize high-dimensional datasets, which are inherently difficult to interpret in their raw form. By projecting data down to two dimensions, users can discover hidden patterns and structures, which is of particular value in exploratory analysis and the presentation of research results.

\section{System Architecture}

The project is designed according to the principles of modularity and separation of concerns, ensuring code readability, ease of extending functionality, and effective management of system complexity.

\subsection{Modular Structure}

The architecture consists of the main application file \texttt{app.py} and six specialized modules located in the \texttt{src/} directory. Each module has a clearly defined responsibility and interface, enabling independent development and testing of individual components.

The \texttt{config.py} module contains the application configuration, method descriptions, and parameter selection guidelines, centralizing all system settings. \texttt{data\_handler.py} is responsible for loading, validating, and preprocessing data, handling both the built-in dataset and user-uploaded files. \texttt{reduction\_methods.py} implements all available dimensionality reduction methods along with their configurable parameters.

Data visualization is handled by the \texttt{visualizations.py} module, which uses the Plotly library to create interactive charts. The \texttt{metrics.py} module provides tools for evaluating the quality of dimensionality reduction by computing various metrics and performing variance analysis. The final module, \texttt{parameter\_optimization.py}, implements parameter optimization through grid search and presents results as interactive heatmaps.

\subsection{Dependency Management and Environment}

The project uses the modern \texttt{uv} package manager, which serves as a fast alternative to pip and poetry. \texttt{uv} provides deterministic package installation, significantly faster operations, and improved virtual environment management. Project configuration is defined in \texttt{pyproject.toml}, enabling standard metadata and dependency management according to PEP 621.

The project uses dependency version locking via \texttt{uv.lock}, guaranteeing environment reproducibility across different machines and development phases. This approach eliminates library version incompatibility issues and ensures application stability.

\subsection{Containerization and Deployment}

The project is additionally prepared for containerization using Docker, ensuring a consistent runtime environment regardless of the target platform. The Dockerfile defines the complete application environment, including all system dependencies and Python configuration.

Through containerization, the application can be easily deployed in various environments --- from local development machines to production servers --- without the need to configure dependencies locally. The container also provides isolation of the application from the host system, enhancing deployment security and stability.

\section{Implemented Dimensionality Reduction Methods}

The application offers six dimensionality reduction methods, each with its own characteristics, configuration parameters, and use cases. The choice of method depends on the nature of the data, structural preservation requirements, and computational constraints.

\subsection{Principal Component Analysis (PCA)}

Principal Component Analysis is a classic linear dimensionality reduction method that finds the directions of maximum variance in data through eigendecomposition of the covariance matrix. PCA is deterministic and computationally efficient, making it ideal for preliminary data analysis and cases requiring fast processing of large datasets.

The method assumes linear relationships between variables and is particularly effective for data with high feature correlation. In the application, PCA also serves as a baseline for comparison with nonlinear methods, allowing users to assess whether the nonlinearity of data justifies the use of more complex algorithms.

\subsection{Kernel Principal Component Analysis (KPCA)}

Kernel PCA extends classical PCA with nonlinear transformations through the use of kernel functions, which map data into a higher-dimensional space where linear relationships can better represent nonlinear structures in the original space.

The application implements four kernel types: RBF (Radial Basis Function), effective for data with local clusters; polynomial kernels, suitable for data with algebraic structure; sigmoidal kernels, mimicking neural network activation functions; and cosine kernels, which exploit angular similarity measures. The gamma parameter controls the influence of individual training points, with higher values leading to more complex decision boundaries.

\subsection{t-Distributed Stochastic Neighbor Embedding (t-SNE)}

t-SNE is a probabilistic nonlinear method particularly effective at visualizing data clusters by preserving local neighborhood structure. The algorithm models similarities between points as probabilities and minimizes the Kullback-Leibler divergence between distributions in the original and reduced spaces.

The key parameter is perplexity, which controls the number of nearest neighbors considered in the analysis. Low perplexity values (5--15) emphasize very local structures, while higher values (30--50) preserve more global relationships. The method is particularly effective at discovering hidden clusters but can be sensitive to parameter selection and computationally expensive for large datasets.

\subsection{Uniform Manifold Approximation and Projection (UMAP)}

UMAP is based on topological manifold theory and offers a trade-off between preserving local and global structure while maintaining high computational efficiency. The algorithm constructs a neighborhood graph in the original space and then optimizes a similar structure in the reduced space.

The \texttt{n\_neighbors} parameter controls the balance between local and global structure --- more neighbors lead to better preservation of global structure at the expense of local detail. The \texttt{min\_dist} parameter specifies the minimum distance between points in the reduced space, where lower values create tighter clusters. UMAP is particularly effective for data with complex topological structure.

\subsection{TriMap}

TRIMAP uses triplet-based relationships between points to preserve both local neighborhoods and global distances. The algorithm defines sets of triplets consisting of a central point along with its nearest and farthest neighbors, and then optimizes the representation in the reduced space to maintain these relationships.

The \texttt{n\_inliers} parameter specifies the number of nearest neighbors considered for each point, affecting the preservation of local structure. The \texttt{n\_outliers} parameter controls the number of distant points used to preserve global structure. TRIMAP often achieves better preservation of global structure than t-SNE with comparable quality of local representation.

In the application, TRIMAP is an optional method: it requires installing an additional dependency group (\texttt{trimap} along with the \texttt{annoy} library, which requires an available C++ compiler). When these dependencies are absent, the method is automatically hidden from the list of available options, and the application informs the user about how to enable it.

\subsection{Pairwise Controlled Manifold Approximation Projection (PaCMAP)}

PaCMAP controls the preservation of local and global structure through precise management of three types of point pairs: near pairs, which preserve local structure; mid-near pairs, which connect local regions; and far pairs, which preserve global structure.

The \texttt{MN\_ratio} parameter controls the proportion of mid-near pairs to near pairs, affecting connections between local regions. The \texttt{FP\_ratio} parameter specifies the ratio of far pairs to near pairs, controlling the preservation of global structure. This precise control makes PaCMAP particularly effective at preserving both local clusters and the global organization of data.

\section{Application Functionality}

The application offers a comprehensive and intuitive workflow that guides the user through all stages of data analysis --- from data loading, through parameter configuration, to detailed results analysis. The interface is designed to be accessible to both novice users and experts requiring advanced configuration options.

\subsection{Data Management and Preprocessing}

The system supports two primary data sources: a built-in diabetes dataset, which serves as a demonstration example, and the ability to upload custom files in CSV and Excel formats. For uploaded files, the application automatically detects column types and allows the user to select a target variable and features for analysis.

The preprocessing module implements a comprehensive data preparation pipeline. The system detects missing values in both features and the target column, after which the user decides whether to drop incomplete rows or halt processing for manual data cleaning. All numeric data are standardized using StandardScaler, which is essential for the correct operation of distance-based methods.

For large datasets, the application offers a sampling feature that limits the number of observations to accelerate computation. The user can specify the number of samples, and the system applies random sampling with a fixed pseudorandom number generator seed, ensuring result reproducibility.

\subsection{Configuration Interface and Parameter Control}

The application sidebar contains dynamically generated parameter controls tailored to the selected dimensionality reduction method. Each method has a defined set of key parameters with appropriate value ranges and default values selected based on best practices.

The system implements intelligent performance warnings that inform users about potentially long computation times for expensive methods (t-SNE, TRIMAP) applied to large datasets. The warnings contain specific suggestions for alternative approaches, such as data sampling or choosing faster methods.

The application also provides detailed parameter selection guidelines for each method, explaining the impact of individual parameters on results and providing recommendations for different data types and analysis goals.

\subsection{Visualization and Results Presentation}

The primary results presentation element is interactive scatter plots created using the Plotly library. Points are colored according to the target variable, enabling immediate assessment of class separation quality in the reduced space. The charts offer full interactivity, including zooming, panning, and displaying detailed information about individual points.

The system automatically adjusts the color scale to the type of target variable --- for categorical variables it uses a discrete color palette, while for continuous variables it applies a color gradient. The charts also include appropriate legends and axis labels.

\subsection{Quality Analysis and Metrics}

The application implements a comprehensive dimensionality reduction quality evaluation system. For linear methods such as PCA, explained variance analysis is computed, showing how much of the total data variance is preserved in the two-dimensional representation.

The system also computes reconstruction metrics for PCA, which measure how well the original data can be reconstructed from the reduced representation. This includes the mean squared reconstruction error and per-feature error distribution visualizations, which help identify variables that are particularly difficult to represent in lower dimensions.

For nonlinear methods, the application displays a limited set of metrics, informing the user about the nature of the selected method. Embedding quality is evaluated using the trustworthiness measure, which quantifies the degree to which neighborhoods in the reduced space correspond to neighborhoods in the original space; this measure is also used during parameter optimization for t-SNE and UMAP.

\subsection{Parameter Optimization}

Advanced parameter optimization functionality enables automatic exploration of the parameter space to find optimal settings for a given dataset. The system implements grid search for key parameters of selected methods: number of components and data subset ratio for PCA, kernel type and gamma for Kernel PCA, perplexity and learning rate for t-SNE, and n\_neighbors and min\_dist for UMAP. Computation results are cached, so repeated analysis runs do not require recomputing already evaluated parameter combinations.

Optimization results are presented as interactive heatmaps that visualize result quality for different parameter combinations. The heatmaps use quality metrics matched to each method: reconstruction error (MSE) for PCA, component variance for Kernel PCA, and trustworthiness for t-SNE and UMAP, allowing the user to select parameters optimal for the specific analysis goal.

The system automatically identifies and highlights the best-performing parameter combinations, and also provides detailed explanations regarding the interpretation of optimization results.

\section{Technical Aspects}

The project uses modern tools for dependency management, in particular the \texttt{uv} package manager, which provides fast installation and a deterministic environment. The application has been containerized with Docker, facilitating deployment and ensuring runtime environment consistency.

Key libraries include Streamlit as the web framework, scikit-learn for classical machine learning methods, and specialized implementations of UMAP, TRIMAP (as an optional dependency), and PaCMAP. Visualizations are realized using Plotly, ensuring interactivity and high graphical quality.

The application implements performance optimization mechanisms, including warnings for large datasets, a sampling feature, and intelligent default parameter selection. Error handling and data validation ensure stable operation across various usage scenarios.

\section{Conclusions and Future Development}

StreamlitDRV constitutes a comprehensive tool for data exploration through dimensionality reduction, combining ease of use with advanced analytical capabilities. The modularity of the architecture enables straightforward extension of functionality with new methods or evaluation metrics.

The project can be developed further in the direction of implementing additional dimensionality reduction methods, expanding statistical analysis capabilities, and integrating with external data sources. Potential improvements also include implementing reduction to more than two dimensions and the ability to export results in various formats.

The application may find use in education, scientific research, and business data analysis, offering an accessible tool for the visualization and exploration of high-dimensional datasets.

\end{document}
